\bilingualchapterstar{Acknowledgements}{致谢}

\begin{englishblock} This book only exists because I spent seven years at AHL learning the craft of systematic trading, although there are no corporate secrets revealed here (as their lawyers will be pleased to know). I wouldn't have been able to write this without spending most of a decade immersed amongst a group of people who almost uniquely in the finance industry managed to be incredibly clever and successful, whilst still being fantastically nice and interesting people to work with. \end{englishblock}

\begin{chineseblock}这本书之所以能够存在，只是因为我曾在 AHL 花了七年时间学习系统化交易这门手艺，尽管这里并没有披露任何公司机密（这一点他们的律师听了应该会很高兴）。如果不是我用了将近十年的时间，沉浸在这样一群人中间，我根本不可能写出这本书。在金融行业里，他们几乎是独一无二的一群人：既极其聪明、极其成功，同时又是非常友善、非常有趣的合作对象。\end{chineseblock}

\bipairgap

\begin{englishblock} There are many people I'd like to thank individually who worked, or still work, at AHL, their parent company the Man Group, and their research centre the Oxford Man Institute; but it would take another book to do it. You know who you are. \end{englishblock}

\begin{chineseblock}我想逐一感谢许多曾经或仍然在 AHL、其母公司英仕曼集团（Man Group）以及研究中心 Oxford Man Institute 工作的人；但真要把每个人都写出来，恐怕还得再写一本书。你们知道自己是谁。\end{chineseblock}

\bipairgap

\begin{englishblock} I'd like to thank my friend Pietro Parodi for his feedback on early drafts of the theoretical chapters where his outside expert viewpoint was invaluable. One former AHL colleague I will mention by name is Thomas Smith. Whilst launching his own hedge fund, Thomas still found the time to spend dozens of unpaid hours reviewing drafts of my work, and behaved exactly like he did when we worked together; he was never afraid of telling me when I had written something that was unintelligible, meaningless, pointless or wrong (and sometimes all of those at the same time). \end{englishblock}

\begin{chineseblock}我要感谢我的朋友 Pietro Parodi，他对理论章节早期草稿提出了反馈，而他那种来自外部专家的视角极其宝贵。还有一位前 AHL 同事我想点名致谢，那就是 Thomas Smith。在筹备他自己的对冲基金期间，Thomas 仍然抽出时间，花了几十个小时无偿审阅我的草稿，而且表现得和我们一起工作时一模一样：只要我写出了难以理解、毫无意义、没有价值或干脆错误的内容，他从来都不会害怕直接告诉我（有时候，一段文字甚至会同时具备以上所有问题）。\end{chineseblock}

\bipairgap

\begin{englishblock} I'm very thankful to Stephen Eckett at Harriman House for discerning the seed of an interesting idea in amongst the turgid ramblings of my original proposal. My editor Craig Pearce has managed the difficult job of reassuring me that I've written a good book, whilst also politely suggesting how it could be improved. \end{englishblock}

\begin{chineseblock}我非常感谢 Harriman House 的 Stephen Eckett，是他从我最初那份冗长乏味、絮絮叨叨的提案中，辨认出了一个有趣想法的种子。我的编辑 Craig Pearce 则完成了一项非常困难的工作：一方面安抚我，让我相信自己写的是一本好书；另一方面又能礼貌地指出，这本书还可以如何改进。\end{chineseblock}

\bipairgap

\begin{englishblock} Last, but certainly not least, I'd like to thank my family. One day they came home and found that their normal office working father and husband had been replaced with a stay at home writer. Worse still he was not creating something interesting like the next Harry Potter, but was instead writing ``Some boring book about money''. Despite that you have still been the most supportive and wonderful people I could ever ask for. Thank you for that, and for everything. \end{englishblock}

\begin{chineseblock}最后，同样也是最重要的，我想感谢我的家人。有一天，他们回到家，发现那个原本正常上班的父亲和丈夫，已经变成了一个待在家里写作的人。更糟的是，他并不是在创作什么像下一部《哈利·波特》那样有趣的东西，而是在写“一本无聊的金钱书”。尽管如此，你们仍然是我所能奢求到的、最支持我、最美好的一群人。感谢你们，感谢这一切。\end{chineseblock}
